% ============================================================
% 03-bezier-matrix-pedagogical-output-explained.tex
% Compile with:  pdflatex -shell-escape 03-bezier-matrix-pedagogical-output-explained.tex
% Run twice for table of contents.
% ============================================================

\documentclass[12pt, a4paper]{article}

% ---- Encoding & fonts ----------------------------------------
\usepackage[T1]{fontenc}
\usepackage[utf8]{inputenc}
\usepackage{lmodern}

% ---- Page geometry -------------------------------------------
\usepackage[
  top=2.5cm, bottom=2.5cm,
  left=2.8cm, right=2.8cm
]{geometry}

% ---- Colours -------------------------------------------------
\usepackage[dvipsnames]{xcolor}
\definecolor{codebg}{RGB}{248,248,248}
\definecolor{rulecolor}{RGB}{180,180,180}
\definecolor{examplebg}{RGB}{240,248,255}
\definecolor{examplerule}{RGB}{100,149,237}
\definecolor{notebg}{RGB}{255,250,230}
\definecolor{noterule}{RGB}{210,180,50}

% ---- Minted --------------------------------------------------
\usepackage{minted}
\setminted[text]{
  bgcolor    = codebg,
  frame      = leftline,
  framerule  = 1.5pt,
  rulecolor  = rulecolor,
  fontsize   = \small,
  breaklines = true,
}

% ---- Math ----------------------------------------------------
\usepackage{amsmath}
\usepackage{amssymb}

% ---- Boxes ---------------------------------------------------
\usepackage{mdframed}

\newmdenv[
  backgroundcolor  = examplebg,
  linecolor        = examplerule, linewidth=1.5pt,
  leftline=true, rightline=false, topline=false, bottomline=false,
  innerleftmargin=10pt, innerrightmargin=8pt,
  innertopmargin=6pt, innerbottommargin=6pt,
  skipabove=6pt, skipbelow=4pt,
]{examplebox}

\newmdenv[
  backgroundcolor  = notebg,
  linecolor        = noterule, linewidth=1.5pt,
  leftline=true, rightline=false, topline=false, bottomline=false,
  innerleftmargin=10pt, innerrightmargin=8pt,
  innertopmargin=6pt, innerbottommargin=6pt,
  skipabove=6pt, skipbelow=4pt,
]{notebox}

% ---- Hyperlinks ----------------------------------------------
\usepackage[colorlinks=true, linkcolor=black, urlcolor=MidnightBlue]{hyperref}

% ---- Misc ----------------------------------------------------
\usepackage{booktabs}
\usepackage{needspace}
\usepackage{parskip}
\usepackage{titlesec}
\usepackage{fancyhdr}
\usepackage{datetime2}
\usepackage{array}

% ---- Section formatting --------------------------------------
\titleformat{\section}
  {\large\bfseries\color{MidnightBlue}}
  {\thesection.}{0.6em}{}[\vspace{-4pt}\rule{\linewidth}{0.4pt}]
\titleformat{\subsection}
  {\normalsize\bfseries\color{BrickRed}}
  {\thesubsection.}{0.5em}{}

% ---- Header / footer -----------------------------------------
\pagestyle{fancy}
\fancyhf{}
\rhead{\textcolor{gray}{\small 03-bezier-matrix-pedagogical-output-explained}}
\lhead{\textcolor{gray}{\small B\'ezier Matrix Method --- Output Explanation}}
\cfoot{\thepage}
\renewcommand{\headrulewidth}{0.3pt}

% ---- Helpers -------------------------------------------------
\newcommand{\cd}[1]{\texttt{#1}}

% ---- Title ---------------------------------------------------
\title{%
  \vspace{1cm}
  {\LARGE\bfseries B\'ezier Curve --- Matrix Method}\\[0.4em]
  {\large $P(u) = U \cdot M \cdot G$}\\[0.4em]
  {\large Explanation of the Program Output}\\[0.4em]
  {\normalsize Case 2: Quartic 3-D, 5 Control Points}\\[0.4em]
  {\normalsize\texttt{03-bezier-matrix-pedagogical.jl}}
}
\author{E.R.\ Nurwijayadi}
\date{\DTMtoday}


% =============================================================
\begin{document}
% =============================================================

\maketitle
\thispagestyle{empty}

\begin{center}
\textit{Directed and curated by E.R.\ Nurwijayadi.
Code and explanations AI-generated.}
\end{center}

\begin{center}
\begin{minipage}{0.85\textwidth}
\small
This document explains the output of
\cd{03-bezier-matrix-pedagogical.jl} when run with \cd{CASE = 2}.
The program computes the same quartic 3-D B\'ezier curve
as the previous scripts, but using an entirely different
mathematical method: the \textbf{matrix method}.
Instead of expanding Bernstein polynomials algebraically,
the curve is expressed as a single matrix product
$P(u) = U \cdot M \cdot G$, which is then evaluated
by substituting powers of $u$ into a row vector.
\end{minipage}
\end{center}

\vspace{1em}\hrule\vspace{1.5em}
\tableofcontents
\newpage


% =============================================================
\needspace{8\baselineskip}
\section{The Big Idea: One Formula, Three Matrices}
% =============================================================

\textit{All previous scripts expanded Bernstein polynomials
term by term.
This script takes a different route:
it packages the entire curve into a single matrix equation
and evaluates it with one matrix multiplication.}

The entire B\'ezier curve is encoded in three matrices:

\[
  P(u) = U(u) \cdot M \cdot G
\]

\begin{center}
\begin{tabular}{clp{7cm}}
\toprule
Matrix & Shape & Role \\
\midrule
$U(u)$ & $1 \times 5$ &
  \textbf{Parameter vector.}
  A row of powers of $u$: $[u^4,\; u^3,\; u^2,\; u,\; 1]$.
  This is the only part that changes when you evaluate at
  different $u$ values. \\[6pt]
$M$ & $5 \times 5$ &
  \textbf{Basis matrix.}
  Encodes all the Bernstein expansion coefficients.
  Computed once from the degree and never changes. \\[6pt]
$G$ & $5 \times 3$ &
  \textbf{Geometry matrix.}
  The control point coordinates, one row per point,
  one column per axis ($x$, $y$, $z$). \\
\bottomrule
\end{tabular}
\end{center}

\medskip
Because $M$ and $G$ do not depend on $u$, the program
computes $M \cdot G$ \emph{once} and stores it.
Evaluating the curve at any $u$ then only requires
building $U(u)$ and computing one dot product.

\begin{notebox}
\textbf{Why this matters.}
With five evaluation points and three axes,
the na\"ive approach calls \cd{evaluate()} fifteen times
(five $u$ values $\times$ three axes).
The matrix method batches all evaluations into a single
matrix multiplication $U_{\text{all}} \cdot (M \cdot G)$
where $U_{\text{all}}$ is a $5 \times 5$ matrix stacking
all $U(u)$ rows.
For large numbers of $u$ values this is dramatically faster.
\end{notebox}


% =============================================================
\needspace{8\baselineskip}
\section{Header and Control Points}
% =============================================================

\begin{minted}{text}
==================================================================
  Bezier Curve - Matrix Method  [P(u) = U * M * G]
  Case 2  |  Quartic  |  3D  |  5 Control Points
==================================================================

  Control Points:
    P0 :  x = 0  y = 0  z = 1
    P1 :  x = 0  y = 4  z = 1
    P2 :  x = 4  y = 0  z = 1
    P3 :  x = 4  y = 4  z = 1
    P4 :  x = 5  y = 4  z = 1
\end{minted}

The subtitle \cd{[P(u) = U * M * G]} states the method upfront.
The control points are the same five points used in the
previous scripts.
They will form the geometry matrix $G$.


% =============================================================
\needspace{8\baselineskip}
\section{Step 1 --- Matrix Formulation}
% =============================================================

\begin{minted}{text}
  STEP 1 - Matrix Formulation:  P(u) = U * M * G
  ----------------------------------------------------------------

  P(u)  =  U(u)  *  M  *  G

  U(u)  shape (1x5)  parameter vector  [u^4, ..., u, 1]
  M     shape (5x5)  basis matrix  (derived from Bernstein)
  G     shape (5x3)  geometry matrix  (control points)

      x  y  z
      -  -  -
  P0  0  0  1
  P1  0  4  1
  P2  4  0  1
  P3  4  4  1
  P4  5  4  1
\end{minted}

This step defines the three matrices and displays $G$ ---
the geometry matrix --- directly as a table of control points.

The geometry matrix $G$ is simply the control points
written as rows:

\[
  G = \begin{bmatrix}
    0 & 0 & 1 \\
    0 & 4 & 1 \\
    4 & 0 & 1 \\
    4 & 4 & 1 \\
    5 & 4 & 1
  \end{bmatrix}
\]

Each row is one control point; each column is one axis.
The table in the output is exactly this matrix, labelled
with axis names and point names.

\begin{notebox}
\textbf{Row and column ordering.}
Rows go $P_0, P_1, \ldots, P_4$ (top to bottom).
Columns go $x, y, z$ (left to right).
This ordering is consistent throughout the output:
$G[i, j]$ is the $j$-th coordinate of the $i$-th control point.
\end{notebox}


% =============================================================
\needspace{8\baselineskip}
\section{Step 2 --- Deriving the Basis Matrix $M$}
% =============================================================

\textit{This is the most mathematical step.
It is also the heart of the matrix method.
Take your time here --- once you understand $M$,
the rest is just arithmetic.}

\needspace{5\baselineskip}
\subsection{What $M$ encodes}

\begin{minted}{text}
  Each column i of M holds the monomial coefficients
  of the Bernstein basis polynomial B(4, i, u):

    B(n,i,u)  =  C(n,i) * u^i * (1-u)^(n-i)

  Expanding (1-u)^(n-i) via the binomial theorem:
    (1-u)^m  =  sum  C(m,j) * (-1)^j * u^j

  Coefficient of u^k in B(n,i,u):
    M[n-k, i]  =  C(n,i) * C(n-i, k-i) * (-1)^(k-i)   for k >= i, else 0

  Row ordering: row 0 = u^4 (highest),  row 4 = u^0 (constant)
\end{minted}

Each column $i$ of $M$ contains the monomial coefficients
of $B(4, i, u)$, reading from the highest power down to the constant.

The row index encodes the power of $u$:
row~0 holds the coefficient of $u^4$,
row~1 holds the coefficient of $u^3$, \ldots,
row~4 holds the constant term ($u^0$).

\begin{examplebox}
\textbf{Why store coefficients this way?}

When we compute $U \cdot M$, each element of the result is
a dot product of the row vector $[u^4, u^3, u^2, u, 1]$
with one column of $M$.
If column $i$ of $M$ holds the coefficients of $B(4,i,u)$
from highest to lowest power, then the dot product
$U \cdot M_{\text{col } i}$ automatically evaluates
$B(4,i,u)$ at the given $u$ --- because it computes
exactly $u^4 \cdot c_4 + u^3 \cdot c_3 + u^2 \cdot c_2
+ u \cdot c_1 + 1 \cdot c_0$.
\end{examplebox}

\needspace{5\baselineskip}
\subsection{The formula for each entry}

The formula shown in the output:
\[
  M[n-k,\; i] = \binom{n}{i} \binom{n-i}{k-i} (-1)^{k-i}
  \quad \text{for } k \geq i, \text{ else } 0
\]

This comes directly from expanding
$B(n,i,u) = \binom{n}{i} u^i (1-u)^{n-i}$
using the binomial theorem on $(1-u)^{n-i}$
and collecting the coefficient of $u^k$.

\needspace{5\baselineskip}
\subsection{Column-by-column derivation}

\begin{minted}{text}
  Column 0  ->  B(4,0,u)  =  (1-u)^4
    M[0,0]  =     1   (coeff of u^4)
    M[1,0]  =    -4   (coeff of u^3)
    M[2,0]  =     6   (coeff of u^2)
    M[3,0]  =    -4   (coeff of u)
    M[4,0]  =     1   (coeff of 1)

  Column 1  ->  B(4,1,u)  =  4u(1-u)^3
    M[0,1]  =    -4   (coeff of u^4)
    M[1,1]  =    12   (coeff of u^3)
    M[2,1]  =   -12   (coeff of u^2)
    M[3,1]  =     4   (coeff of u)
    ...
\end{minted}

Each column is shown with its non-zero entries only.
Let us verify column~0 by hand.

\begin{examplebox}
\textbf{Verify column~0: $B(4,0,u) = (1-u)^4$}

Expand $(1-u)^4$ using the binomial theorem:
\[
  (1-u)^4 = \sum_{j=0}^{4} \binom{4}{j}(-1)^j u^j
  = 1 - 4u + 6u^2 - 4u^3 + u^4
\]

Reading coefficients from highest power down:
\[
  u^4 \to +1, \quad u^3 \to -4, \quad u^2 \to +6,
  \quad u^1 \to -4, \quad u^0 \to +1
\]

These are exactly $M[0,0]=1$, $M[1,0]=-4$, $M[2,0]=6$,
$M[3,0]=-4$, $M[4,0]=1$. \checkmark

\medskip
\textbf{Verify column~2: $B(4,2,u) = 6u^2(1-u)^2$}

$(1-u)^2 = 1 - 2u + u^2$, so:
\[
  6u^2(1 - 2u + u^2) = 6u^2 - 12u^3 + 6u^4
\]

Only three non-zero entries (the $u^0$ and $u^1$ terms are zero):
\[
  u^4 \to +6, \quad u^3 \to -12, \quad u^2 \to +6
\]

So $M[0,2]=6$, $M[1,2]=-12$, $M[2,2]=6$, and all lower rows
in column~2 are zero. \checkmark
\end{examplebox}

\needspace{5\baselineskip}
\subsection{The full $M$ matrix}

\begin{minted}{text}
  Full M matrix:

       P0   P1   P2  P3  P4
       --  ---  ---  --  --
  u^4   1   -4    6  -4   1
  u^3  -4   12  -12   4   0
  u^2   6  -12    6   0   0
  u^1  -4    4    0   0   0
  u^0   1    0    0   0   0
\end{minted}

Written as a matrix (rows = powers from $u^4$ down, columns = $B_0$ through $B_4$):

\[
  M = \begin{bmatrix}
    1  & -4 &  6 & -4 & 1 \\
   -4  & 12 & -12 & 4 & 0 \\
    6  & -12 &  6 &  0 & 0 \\
   -4  &  4 &  0 &  0 & 0 \\
    1  &  0 &  0 &  0 & 0
  \end{bmatrix}
\]

\begin{notebox}
\textbf{Pattern in $M$.}
The lower-left triangle is all zeros --- this is because
$B(4,i,u)$ contains $u^i$ as its lowest power, so all
powers below $u^i$ in column $i$ are zero.
The matrix is therefore \textbf{upper-triangular}
(in the sense that all entries below the anti-diagonal are zero).

The first column contains the coefficients of $(1-u)^4$,
which is the familiar $1, -4, 6, -4, 1$ pattern from Pascal's triangle
with alternating signs.
The last column contains a single $1$ in the top row ---
because $B(4,4,u) = u^4$ has only one term.
\end{notebox}


% =============================================================
\needspace{8\baselineskip}
\section{Step 3 --- Computing $M \cdot G$}
% =============================================================

\begin{minted}{text}
  STEP 3 - Compute  M * G  - Polynomial Coefficient Matrix
  ----------------------------------------------------------------

  M * G  gives the monomial coefficients for each axis.
  Row k of (M*G) is the coefficient of u^(n-k).

         x    y  z
       ---  ---  -
  u^4   13  -28  0
  u^3  -32   64  0
  u^2   24  -48  0
  u^1    0   16  0
  u^0    0    0  1

  Resulting polynomials:
    x(u)  =  13u^4  - 32u^3  + 24u^2
    y(u)  =  -28u^4  + 64u^3  - 48u^2  + 16u
    z(u)  =  1
\end{minted}

$M \cdot G$ is a $5 \times 3$ matrix.
Each row gives the coefficient of one power of $u$
for all three axes simultaneously.
Row~0 (labelled $u^4$) gives the leading coefficients;
row~4 (labelled $u^0$) gives the constant terms.

\begin{examplebox}
\textbf{Verify the $x$ column of $M \cdot G$:}

The $x$ column of $G$ is $[0, 0, 4, 4, 5]^T$.
Multiplying $M \cdot [0,0,4,4,5]^T$ row by row:

\medskip
\textbf{Row 0} (coeff of $u^4$):
\[
  1(0) + (-4)(0) + 6(4) + (-4)(4) + 1(5)
  = 0 + 0 + 24 - 16 + 5 = 13 \checkmark
\]

\textbf{Row 1} (coeff of $u^3$):
\[
  (-4)(0) + 12(0) + (-12)(4) + 4(4) + 0(5)
  = 0 + 0 - 48 + 16 + 0 = -32 \checkmark
\]

\textbf{Row 2} (coeff of $u^2$):
\[
  6(0) + (-12)(0) + 6(4) + 0(4) + 0(5)
  = 0 + 0 + 24 + 0 + 0 = 24 \checkmark
\]

\textbf{Rows 3 and 4} (coeff of $u$ and constant):
\[
  (-4)(0) + 4(0) + 0(4) + 0(4) + 0(5) = 0 \checkmark
\]
\[
  1(0) + 0(0) + 0(4) + 0(4) + 0(5) = 0 \checkmark
\]

So the $x$ column of $M \cdot G$ is $[13, -32, 24, 0, 0]^T$,
giving $x(u) = 13u^4 - 32u^3 + 24u^2$. \checkmark
\end{examplebox}

\begin{examplebox}
\textbf{Verify the $z$ column of $M \cdot G$:}

The $z$ column of $G$ is $[1, 1, 1, 1, 1]^T$.
Each row of $M$ sums to zero except the constant row:

The columns of $M$ are the monomial expansions of
$B(4,0)$ through $B(4,4)$.
Their sum is the expansion of $\sum_{i=0}^{4} B(4,i,u) = 1$,
which in monomial form is simply $u^0 = 1$ ---
all higher powers cancel.

Therefore $M \cdot [1,1,1,1,1]^T = [0,0,0,0,1]^T$,
giving $z(u) = 1$. \checkmark
\end{examplebox}


% =============================================================
\needspace{8\baselineskip}
\section{Step 4 --- Evaluating $U \cdot (M \cdot G)$}
% =============================================================

\textit{This step is the most explicit in the entire output.
For each $u$ value, the program shows the full $U$ vector
and then writes out every term of every dot product.
Nothing is hidden.}

\needspace{5\baselineskip}
\subsection{Structure of each block}

Each $u$ value gets its own block with two parts:
the $U$ vector table, then the dot product expansion.

\begin{minted}{text}
  -- u = 0.5 -------------------------------------------------

  U  =
    u^4   u^3   u^2   u^1  u^0
  -----  ----  ----  ----  ---
  1//16  1//8  1//4  1//2    1

  U * (M*G)  =  P(u):
    x(u)  =  1//16*(13)  +  1//8*(-32)  +  1//4*(24)
             +  1//2*(0)  +  1*(0)
           =  45//16
    y(u)  =  ...
\end{minted}

The $U$ vector is displayed as a single-row table
with each power of $u$ in its own column.
The values are shown as exact fractions (\cd{1//16}, \cd{1//8}, etc.)
because the program uses \cd{fmt\_frac} which converts floats
back to their nearest rational representation.

\needspace{5\baselineskip}
\subsection{Reading the $U$ vector}

The $U$ vector for $u = \tfrac{1}{2}$ is:
\[
  U = \left[\left(\tfrac{1}{2}\right)^4,\;
            \left(\tfrac{1}{2}\right)^3,\;
            \left(\tfrac{1}{2}\right)^2,\;
            \tfrac{1}{2},\;
            1\right]
  = \left[\tfrac{1}{16},\; \tfrac{1}{8},\; \tfrac{1}{4},\;
          \tfrac{1}{2},\; 1\right]
\]

These are exactly the values shown in the output table:
\cd{1//16}, \cd{1//8}, \cd{1//4}, \cd{1//2}, \cd{1}.

\needspace{5\baselineskip}
\subsection{Reading the dot product lines}

Each axis line shows:
\[
  \text{axis}(u) = U[0] \cdot (MG)[0] + U[1] \cdot (MG)[1]
  + \cdots + U[4] \cdot (MG)[4]
\]

where $(MG)[k]$ is the entry from the pre-computed $M \cdot G$
column for that axis, at row $k$.

\begin{examplebox}
\textbf{Hand-verify $x\!\left(\tfrac{1}{2}\right) = \dfrac{45}{16}$:}

From Step~3, the $x$ column of $M \cdot G$ is $[13, -32, 24, 0, 0]$.
The $U$ vector at $u = \tfrac{1}{2}$ is
$[\tfrac{1}{16}, \tfrac{1}{8}, \tfrac{1}{4}, \tfrac{1}{2}, 1]$.

Dot product:
\[
  \tfrac{1}{16} \cdot 13
  + \tfrac{1}{8} \cdot (-32)
  + \tfrac{1}{4} \cdot 24
  + \tfrac{1}{2} \cdot 0
  + 1 \cdot 0
\]
\[
  = \tfrac{13}{16} - \tfrac{32}{8} + \tfrac{24}{4}
  = \tfrac{13}{16} - 4 + 6
  = \tfrac{13}{16} + 2
  = \tfrac{13 + 32}{16}
  = \tfrac{45}{16}
\]

Displayed as \cd{45//16} $= 2.8125$. \checkmark
\end{examplebox}

\begin{examplebox}
\textbf{Hand-verify $y\!\left(\tfrac{1}{4}\right) = \dfrac{121}{64}$:}

From Step~3, the $y$ column of $M \cdot G$ is $[-28, 64, -48, 16, 0]$.
The $U$ vector at $u = \tfrac{1}{4}$ is
$[\tfrac{1}{256}, \tfrac{1}{64}, \tfrac{1}{16}, \tfrac{1}{4}, 1]$.

Dot product (converting to denominator 256):
\[
  \tfrac{1}{256}(-28) + \tfrac{1}{64}(64) + \tfrac{1}{16}(-48)
  + \tfrac{1}{4}(16) + 1(0)
\]
\[
  = \tfrac{-28}{256} + \tfrac{256}{256} + \tfrac{-768}{256}
  + \tfrac{1024}{256}
  = \tfrac{-28 + 256 - 768 + 1024}{256}
  = \tfrac{484}{256}
  = \tfrac{121}{64}
\]

Displayed as \cd{121//64} $\approx 1.8906$. \checkmark
\end{examplebox}

\needspace{5\baselineskip}
\subsection{The boundary cases}

\begin{minted}{text}
  -- u = 0.0 -------------------------------------------------
  U  =  [0  0  0  0  1]

  -- u = 1.0 -------------------------------------------------
  U  =  [1  1  1  1  1]
\end{minted}

\begin{examplebox}
\textbf{At $u = 0$:} $U = [0, 0, 0, 0, 1]$.

The dot product with any column of $M \cdot G$ picks
only the last row --- the constant term (row~4, $u^0$).

For $x$: the constant term is $0$. Result: $x(0) = 0$.
For $z$: the constant term is $1$. Result: $z(0) = 1$.

So $P(0) = (0, 0, 1) = P_0$. \checkmark
The curve starts at the first control point.

\medskip
\textbf{At $u = 1$:} $U = [1, 1, 1, 1, 1]$.

The dot product sums all rows of the $M \cdot G$ column.

For $x$: $13 + (-32) + 24 + 0 + 0 = 5$. \checkmark
For $y$: $-28 + 64 + (-48) + 16 + 0 = 4$. \checkmark

So $P(1) = (5, 4, 1) = P_4$. \checkmark
The curve ends at the last control point.
\end{examplebox}


% =============================================================
\needspace{8\baselineskip}
\section{Step 5 --- Summary Table}
% =============================================================

\begin{minted}{text}
  STEP 5 - Summary - Evaluation Table
  ----------------------------------------------------------------

  +=========================================================+
  | u       x(u)               y(u)             z(u)       |
  |=========================================================|
  | 0       0  (0.0000)        0  (0.0000)  1  (1.0000)    |
  | 0.2500  269//256  (1.0508) 121//64 (1.8906) 1 (1.0000) |
  | 0.5000  45//16  (2.8125)   9//4  (2.2500)  1  (1.0000) |
  | 0.7500  1053//256  (4.1133)201//64(3.1406)  1  (1.0000) |
  | 1       5  (5.0000)        4  (4.0000)  1  (1.0000)    |
  +==========================================================+
\end{minted}

The final summary table collects all results.
It uses a box-border style (with Unicode box-drawing characters
\cd{+}, \cd{|}, \cd{=}) to visually distinguish
it from the plain tables used in earlier steps.

Each cell shows the exact rational result followed by the
4-decimal-place float in parentheses:
\cd{269//256~(1.0508)} means the exact value is $\tfrac{269}{256}$
and the decimal approximation is $1.0508$.

\begin{notebox}
\textbf{All results match the previous scripts.}
The five points on the curve are exactly the same as those
computed by the Bernstein expansion in scripts~01 and~02.
This is the expected result: the matrix method and the
algebraic expansion are mathematically equivalent.
They arrive at the same answer by different routes.

$u=0$: $(0, 0, 1) = P_0$. The curve starts at $P_0$. \\
$u=1$: $(5, 4, 1) = P_4$. The curve ends at $P_4$. \\
$z(u) = 1$ for all $u$: the curve lies flat in the plane $z=1$.
\end{notebox}


% =============================================================
\needspace{8\baselineskip}
\section{Why the Matrix Method?}
% =============================================================

\textit{The matrix method produces the same curve.
So why use it at all?
Three reasons.}

\textbf{Efficiency.}
$M \cdot G$ is computed once.
Evaluating the curve at $N$ different $u$ values then
costs only one $N \times (n+1)$ matrix multiplication.
Compare this to expanding and evaluating $n+1$ Bernstein
polynomials separately for each of $N$ values.

\textbf{Insight.}
The matrix $M$ is fixed for any curve of the same degree.
Only $G$ changes when you move control points.
This makes it easy to see how changing one control point
affects the coefficients: you only need to update
one row of $G$ and recompute $M \cdot G$.

\textbf{Transparency.}
Step~4 in this output is the most explicit evaluation
shown in any of the three scripts.
Every multiplication is visible.
There is no recursion, no accumulation ---
just a row vector times a pre-computed matrix.

\begin{examplebox}
\textbf{Summary of the computation chain:}

\medskip
\begin{tabular}{rl}
\toprule
Step & Operation \\
\midrule
Once & $M \gets \text{derive\_M}(\text{degree})$ \\
Once & $MG \gets M \cdot G$ \\
Per $u$ & $U \gets [u^4, u^3, u^2, u, 1]$ \\
Per $u$ & $P(u) \gets U \cdot MG$ \\
\bottomrule
\end{tabular}

\medskip
The heavy work ($M$ derivation and $M \cdot G$ multiplication)
is done once.
All repeated evaluations are cheap.
\end{examplebox}

\vspace{2em}\hrule\vspace{0.5em}
{\small\textcolor{gray}{%
Output explanation for \texttt{03-bezier-matrix-pedagogical.jl}, Case~2
\textbar{} Directed and curated by E.R.\ Nurwijayadi
\textbar{} Code and explanations AI-generated}}

\end{document}
